
\section{Bottom layer: noncommutative scanning and finite-resolution readout}
\label{sec:noncommutative_scanning_readout}

\subsection{The irrational rotation algebra as the minimal scan closure}
\label{sec:rotation_algebra}

Under the Weyl relation of Axiom~O6, the smallest $C^\ast$-algebraic closure generated by the scan/readout unitaries is the \emph{irrational rotation algebra} (the standard noncommutative $2$-torus)
\begin{equation}
\label{eq:rotation_algebra_def}
A_\alpha := C^\ast\langle U,V \mid UV=\e^{2\pi \iu \alpha}VU\rangle,\qquad \alpha\in(0,1)\setminus\QQ,
\end{equation}
studied classically by Rieffel and Connes \cite{Rieffel1981,Connes1994}.
In the protocol semantics, $U$ is a tick shift (scan iteration) and $V$ is a pointer phase unitary, while the noncommutativity in~\eqref{eq:rotation_algebra_def} encodes intrinsic incompatibility between ``scan time'' and ``readout phase''.

\begin{remark}[rigidity of the irrational rotation algebra]
For irrational $\alpha$, the rotation algebra $A_\alpha$ is simple and admits a unique tracial state; it is a canonical, rigid noncommutative geometric object associated with an irrational slope \cite{Rieffel1981,Connes1994}.
\end{remark}

\begin{proposition}[no simultaneous eigenvector]
\label{prop:no_simultaneous_eigenvector}
Let $\alpha\notin\QQ$ and let $U,V$ be unitaries satisfying $UV=\e^{2\pi\iu\alpha}VU$. Then there is no nonzero vector $\psi$ such that $U\psi=\lambda\psi$ and $V\psi=\mu\psi$ for some $\lambda,\mu\in\CC$.
\end{proposition}
\begin{proof}
If such $\psi$ existed, then $UV\psi=\lambda\mu\psi$ and $VU\psi=\mu\lambda\psi$. The Weyl relation yields $\lambda\mu=\e^{2\pi\iu\alpha}\mu\lambda$. Since $|\lambda|=|\mu|=1$ and $\e^{2\pi\iu\alpha}\neq 1$ for $\alpha\notin\QQ$, this is impossible.
\end{proof}

\subsection{A canonical covariant model on the circle}
\label{sec:covariant_model}

A standard realization of~\eqref{eq:rotation_algebra_def} is on $\hilbert=L^2(\RR/\ZZ)$:
\begin{equation}
\label{eq:covariant_model}
(U\psi)(x)=\psi(x+\alpha),\qquad (V\psi)(x)=\e^{2\pi\iu x}\psi(x).
\end{equation}
The induced classical orbit on the circle is the irrational rotation
\begin{equation}
x_t = x_0 + t\alpha \pmod 1,
\end{equation}
which is uniquely ergodic and equidistributed \cite{Denjoy1932,KatokHasselblatt1995,Weyl1916,KuipersNiederreiter1974}.

\subsection{Finite-resolution readout: windows, effects, and induced statistics}
\label{sec:finite_resolution_readout}

Finite-resolution readout is \emph{not} an external sampling axiom; it is part of the protocol.
The simplest example is a window $W\subset\RR/\ZZ$ and a binary readout
\begin{equation}
b_t := \ind_W(x_t)\in\{0,1\}.
\end{equation}
More generally, let $\{w_k^{(\varepsilon)}(x)\}_{k\in K}$ be a measurable partition of unity with $w_k^{(\varepsilon)}(x)\in[0,1]$ and $\sum_k w_k^{(\varepsilon)}(x)=1$. Using the spectral measure $\Pi_V$ of the phase unitary $V$, define effects
\begin{equation}
E_k^{(\varepsilon)} := \int_{\RR/\ZZ} w_k^{(\varepsilon)}(x)\,\dd \Pi_V(x),\qquad \sum_k E_k^{(\varepsilon)}=\ind,
\end{equation}
and probabilities $P_k^{(\varepsilon)}=\omega_{\mathrm{eff}}(E_k^{(\varepsilon)})$ as in O5.

\begin{theorem}[Weyl equidistribution and induced frequencies]
\label{thm:weyl_equidistribution}
If $\alpha\notin\QQ$, the orbit $x_t=x_0+t\alpha\pmod 1$ is equidistributed in $\RR/\ZZ$. In particular, for any interval window $W$,
\begin{equation}
\frac{1}{N}\sum_{t=0}^{N-1}\ind_W(x_t)\;\longrightarrow\; |W|,\qquad N\to\infty,
\end{equation}
where $|W|$ is the Lebesgue length of $W$.
\end{theorem}
\begin{remark}
The statement extends to Riemann-integrable functions and, under bounded-variation assumptions, admits explicit finite-$N$ discrepancy bounds; see Section~\ref{sec:cusp_qexp_slice_sampling}. Further details are in Appendix~\ref{app:dk_koksma_ostrowski}.
\end{remark}

\subsection{Symbolic readout and the golden branch}
\label{sec:sturmian_golden}

When $W$ is an interval, the binary readout $b_t=\ind_W(x_t)$ is a \emph{Sturmian word} (a mechanical word) of minimal complexity $p(n)=n+1$ \cite{MorseHedlund1940,Lothaire2002}.
The ``golden branch'' $\alpha=\varphi^{-1}$ (continued fraction $[0;1,1,1,\dots]$) yields Fibonacci/Sturmian structure and, at the coding interface, the Zeckendorf decomposition. This branch plays the role of a canonical toy sector for explicit bounds (Section~\ref{sec:quantitative_closure}).


